% GENERATED FILE. Edit canonical lessons, then run npm run generate:books.
% canonical-source-hash: fnv1a64:a05e46b49008c403
% canonical-lessons: BN-C295-khoti, BN-C295-dhvongsostup, BN-C295-meramot-kora, BN-C295-mokabila-kora, BN-C295-phire-paoya

\chapter{Counting the Cost}
\label{ch:bn-counting-the-cost}
\hlchaptermodality{295}

\begin{chapteropening}
\textbf{By the end of this chapter:} \emph{I can talk about damage, hardship and getting over a disaster.}

Complete the last lesson of chapter 295: I can talk about damage, hardship and getting over a disaster.
\end{chapteropening}

\section[\texorpdfstring{\bn{ক্ষতি} (khôti)}{ক্ষতি (khôti)}]{\bn{ক্ষতি} (khôti) — damage; harm; a loss}

\label{lesson:BN-C295-khoti}



\begin{quote}
Before the new one: say the Bengali for an apprehension, then the Bengali for to avoid.
\end{quote}

\subsection*{You'll want to know: \bn{ক্ষতি}}
\textbf{\bn{ক্ষতি}} — \emph{khôti} — \textquotedblleft{}damage; harm; a loss\textquotedblright{}.

\begin{quote}
\bn{বন্যায়} \bn{গ্রামের} \bn{অনেক} \bn{ক্ষতি} \bn{হল।}
\end{quote}

\begin{quote}
\emph{bônyāy grāmer ônek khôti hôlô.}
\end{quote}

\begin{quote}
The flood did a lot of damage to the village.
\end{quote}

\begin{cousinweb}
\textbf{\bn{ক্ষতি}} (\emph{khôti}) is Sanskrit \emph{kṣati}, 'injury, hurt; loss' (Monier-Williams). \bn{ক্ষতি} \bn{হওয়া}, to suffer damage; \bn{ক্ষতি} \bn{করা}, to do harm.
\end{cousinweb}

\begin{grammarlens}[title={What you've built}]
One more word for telling what happened.
\end{grammarlens}

\subsection*{Your turn}
\begin{itemize}
\raggedright
  \item \emph{Say it:} \emph{khôti}
  \item \emph{Say it:} \emph{khôti}, once more
  \item \emph{Say it:} \emph{bônyāy grāmer ônek khôti hôlô.} — The flood did a lot of damage to the village.
  \item \emph{Recall:} say the Bengali for an apprehension, then the Bengali for to avoid, then say \emph{khôti} again
\end{itemize}

\begin{tcolorbox}[breakable,colback=teal!4,colframe=teal!35!black,title={Before you move on}]
What does \bn{ক্ষতি} mean? (\textquotedblleft{}Damage; harm; a loss\textquotedblright{}.) Say it once more.
\end{tcolorbox}
\section[\texorpdfstring{\bn{ধ্বংসস্তূপ} (dhvôngsôstūp)}{ধ্বংসস্তূপ (dhvôngsôstūp)}]{\bn{ধ্বংসস্তূপ} (dhvôngsôstūp) — rubble; a heap of ruins}

\label{lesson:BN-C295-dhvongsostup}



\begin{quote}
Before the new one: say the Bengali for to avoid, then the Bengali for damage.
\end{quote}

\subsection*{You'll want to know: \bn{ধ্বংসস্তূপ}}
\textbf{\bn{ধ্বংসস্তূপ}} — \emph{dhvôngsôstūp} — \textquotedblleft{}rubble; a heap of ruins\textquotedblright{}.

\begin{quote}
\bn{স্বেচ্ছাসেবকরা} \bn{ধ্বংসস্তূপ} \bn{সরিয়ে} \bn{রাস্তা} \bn{খুলল।}
\end{quote}

\begin{quote}
\emph{svechchhāsebôkrā dhvôngsôstūp sôriye rāstā khullô.}
\end{quote}

\begin{quote}
The volunteers cleared the rubble and opened the road.
\end{quote}

\begin{cousinweb}
\textbf{\bn{ধ্বংসস্তূপ}} (\emph{dhvôngsôstūp}) joins \textbf{\bn{ধ্বংস}}, ruin, to Sanskrit \emph{stūpa}, 'a heap, a mound' (Monier-Williams), the word behind the Buddhist stupa: a heap of what was destroyed.
\end{cousinweb}

\begin{grammarlens}[title={What you've built}]
One more word for telling what happened.
\end{grammarlens}

\subsection*{Your turn}
\begin{itemize}
\raggedright
  \item \emph{Say it:} \emph{dhvôngsôstūp}
  \item \emph{Say it:} \emph{dhvôngsôstūp}, once more
  \item \emph{Say it:} \emph{svechchhāsebôkrā dhvôngsôstūp sôriye rāstā khullô.} — The volunteers cleared the rubble and opened the road.
  \item \emph{Recall:} say the Bengali for to avoid, then the Bengali for damage, then say \emph{dhvôngsôstūp} again
\end{itemize}

\begin{tcolorbox}[breakable,colback=teal!4,colframe=teal!35!black,title={Before you move on}]
What does \bn{ধ্বংসস্তূপ} mean? (\textquotedblleft{}Rubble; a heap of ruins\textquotedblright{}.) Say it once more.
\end{tcolorbox}
\section[\texorpdfstring{\bn{মেরামত} \bn{করা} (merāmôt kôrā)}{মেরামত করা (merāmôt kôrā)}]{\bn{মেরামত} \bn{করা} (merāmôt kôrā) — to mend; to patch up}

\label{lesson:BN-C295-meramot-kora}



\begin{quote}
Before the new one: say the Bengali for damage, then the Bengali for rubble.
\end{quote}

\subsection*{You'll want to know: \bn{মেরামত} \bn{করা}}
\textbf{\bn{মেরামত} \bn{করা}} — \emph{merāmôt kôrā} — \textquotedblleft{}to mend; to patch up\textquotedblright{}.

\begin{quote}
\bn{স্বেচ্ছাসেবকরা} \bn{ভাঙা} \bn{সেতুটা} \bn{মেরামত} \bn{করল।}
\end{quote}

\begin{quote}
\emph{svechchhāsebôkrā bhāngā setuṭā merāmôt kôrlô.}
\end{quote}

\begin{quote}
The volunteers patched up the broken bridge.
\end{quote}

Past (he/she …): \bn{তিনি} \bn{ছাদটা} \bn{মেরামত} \bn{করলেন}, \emph{tini chhādṭā merāmôt kôrlen}.

\begin{cousinweb}
\textbf{\bn{মেরামত} \bn{করা}} (\emph{merāmôt kôrā}) puts \bn{মেরামত} — usually traced through Persian to Arabic \emph{marammat}, 'repair' — before \bn{করা}. It is the word for fixing what a storm broke; \textbf{\bn{সারানো}} is the everyday word for repairing.
\end{cousinweb}

\begin{grammarlens}[title={What you've built}]
One more word for telling what happened.
\end{grammarlens}

\subsection*{Your turn}
\begin{itemize}
\raggedright
  \item \emph{Say it:} \emph{merāmôt kôrā}
  \item \emph{Say it:} \emph{merāmôt kôrā}, once more
  \item \emph{Say it:} \emph{svechchhāsebôkrā bhāngā setuṭā merāmôt kôrlô.} — The volunteers patched up the broken bridge.
  \item \emph{Recall:} say the Bengali for damage, then the Bengali for rubble, then say \emph{merāmôt kôrā} again
\end{itemize}

\begin{tcolorbox}[breakable,colback=teal!4,colframe=teal!35!black,title={Before you move on}]
What does \bn{মেরামত} \bn{করা} mean? (\textquotedblleft{}To mend; to patch up\textquotedblright{}.) Say it once more.
\end{tcolorbox}
\section[\texorpdfstring{\bn{মোকাবিলা} \bn{করা} (mokābilā kôrā)}{মোকাবিলা করা (mokābilā kôrā)}]{\bn{মোকাবিলা} \bn{করা} (mokābilā kôrā) — to face; to tackle}

\label{lesson:BN-C295-mokabila-kora}



\begin{quote}
Before the new one: say the Bengali for rubble, then the Bengali for to mend.
\end{quote}

\subsection*{You'll want to know: \bn{মোকাবিলা} \bn{করা}}
\textbf{\bn{মোকাবিলা} \bn{করা}} — \emph{mokābilā kôrā} — \textquotedblleft{}to face; to tackle\textquotedblright{}.

\begin{quote}
\bn{গ্রামের} \bn{মানুষ} \bn{সাহস} \bn{নিয়ে} \bn{বন্যার} \bn{মোকাবিলা} \bn{করল।}
\end{quote}

\begin{quote}
\emph{grāmer mānush sāhôs niye bônyār mokābilā kôrlô.}
\end{quote}

\begin{quote}
The villagers faced the flood with courage.
\end{quote}

Past (he/she …): \bn{তিনি} \bn{বিপদের} \bn{মোকাবিলা} \bn{করলেন}, \emph{tini bipôder mokābilā kôrlen}.

\begin{cousinweb}
\textbf{\bn{মোকাবিলা} \bn{করা}} (\emph{mokābilā kôrā}) puts \bn{মোকাবিলা} — usually traced through Persian to Arabic \emph{muqābala}, 'facing each other' — before \bn{করা}: to stand face to face with a trouble.
\end{cousinweb}

\begin{grammarlens}[title={What you've built}]
One more word for telling what happened.
\end{grammarlens}

\subsection*{Your turn}
\begin{itemize}
\raggedright
  \item \emph{Say it:} \emph{mokābilā kôrā}
  \item \emph{Say it:} \emph{mokābilā kôrā}, once more
  \item \emph{Say it:} \emph{grāmer mānush sāhôs niye bônyār mokābilā kôrlô.} — The villagers faced the flood with courage.
  \item \emph{Recall:} say the Bengali for rubble, then the Bengali for to mend, then say \emph{mokābilā kôrā} again
\end{itemize}

\begin{tcolorbox}[breakable,colback=teal!4,colframe=teal!35!black,title={Before you move on}]
What does \bn{মোকাবিলা} \bn{করা} mean? (\textquotedblleft{}To face; to tackle\textquotedblright{}.) Say it once more.
\end{tcolorbox}
\section[\texorpdfstring{\bn{ফিরে} \bn{পাওয়া} (phire pāoyā)}{ফিরে পাওয়া (phire pāoyā)}]{\bn{ফিরে} \bn{পাওয়া} (phire pāoyā) — to get back; to regain}

\label{lesson:BN-C295-phire-paoya}



\begin{quote}
Before the new one: say the Bengali for to mend, then the Bengali for to tackle.
\end{quote}

\subsection*{You'll want to know: \bn{ফিরে} \bn{পাওয়া}}
\textbf{\bn{ফিরে} \bn{পাওয়া}} — \emph{phire pāoyā} — \textquotedblleft{}to get back; to regain\textquotedblright{}.

\begin{quote}
\bn{অনেক} \bn{মাস} \bn{পরে} \bn{গ্রামটা} \bn{আবার} \bn{তার} \bn{পুরনো} \bn{চেহারা} \bn{ফিরে} \bn{পেল।}
\end{quote}

\begin{quote}
\emph{ônek mās pôre grāmṭā ābār tār purôno chehārā phire pelô.}
\end{quote}

\begin{quote}
Months later the village got its old look back.
\end{quote}

Past (he/she …): \bn{সেই} \bn{লোক} \bn{তার} \bn{ব্যাগ} \bn{ফিরে} \bn{পেল}, \emph{sei lok tār byāg phire pelô}.

\begin{cousinweb}
\textbf{\bn{ফিরে} \bn{পাওয়া}} (\emph{phire pāoyā}) is the -e form of \textbf{\bn{ফেরা}}, to come back — CDIAL 9078, a Middle Indo-Aryan \emph{phirati}, 'turns' — before \textbf{\bn{পাওয়া}}, to get: to get something back, as \textbf{\bn{খুঁজে} \bn{পাওয়া}} is to find by searching.
\end{cousinweb}

\begin{grammarlens}[title={What you've built}]
One more word for telling what happened.
\end{grammarlens}

\subsection*{Your turn}
\begin{itemize}
\raggedright
  \item \emph{Say it:} \emph{phire pāoyā}
  \item \emph{Say it:} \emph{phire pāoyā}, once more
  \item \emph{Say it:} \emph{ônek mās pôre grāmṭā ābār tār purôno chehārā phire pelô.} — Months later the village got its old look back.
  \item \emph{Recall:} say the Bengali for to mend, then the Bengali for to tackle, then say \emph{phire pāoyā} again
\end{itemize}

\begin{tcolorbox}[breakable,colback=teal!4,colframe=teal!35!black,title={Before you move on}]
What does \bn{ফিরে} \bn{পাওয়া} mean? (\textquotedblleft{}To get back; to regain\textquotedblright{}.) Say it once more.
\end{tcolorbox}
